% Source: https://artofproblemsolving.com/community/c3h2328545_beautiful_random_problems
% Stable identity: aops-topic-2328545-q01
Each day the remaining amount is multiplied by $1-\dfrac16=\dfrac56$, so after $n$ days $\left(\dfrac56\right)^n$ of the cake is left.

\textbf{(a)} $\left(\dfrac56\right)^n>0$ for every $n$, so some cake is always left: I \textbf{never} finish the cake (if it can be divided indefinitely). The remaining amount does shrink towards $0$.

\textbf{(b)} $\left(\dfrac56\right)^3=\dfrac{125}{216}>\dfrac12$, while $\left(\dfrac56\right)^4=\dfrac{625}{1296}<\dfrac12$ (since $2\cdot625=1250<1296$). So less than half is left for the first time after $\boxed{4}$ days.
